% ABSTRACT (English)
% Note: Should not exceed 300 words.
% Should cover: Objective, Methodology, Results, Applications
% Use Past tense for objectives and methodology
% Use Present tense for results and applications
%
% *** Delete the example text below and write your own abstract ***

This project presented an application of the end-to-end data science process to analyze key pre-match factors in European football leagues and to develop predictive models for predicting three-way match outcomes (Home Win, Draw, Away Win) to identify positive expected value (+EV) betting opportunities and evaluate financial returns.

The methodology integrated historical match data across five major European leagues from the 2014--2015 through 2024--2025 seasons from public sources, including Understat, Transfermarkt, and Football-Data.co.uk. Data was processed through a modular data pipeline enforced with Pandera schema validation. Pre-match feature engineering was performed to construct contextual metrics, including Elo ratings, cumulative fatigue indices, schedule congestion across all competitions, and starting lineup forms (xG/xA per 90 minutes). A Dixon-Coles adjusted Double Poisson statistical model and a HistGradientBoosting classifier were trained, with financial viability evaluated via a temporal backtesting engine using Fixed Staking and Kelly Criterion strategies.

The experimental results indicate that integrating pre-match contextual features—specifically Elo ratings, fatigue metrics, and starting lineup forms—with Gradient Boosting and Double Poisson models significantly enhances outcome probability estimation beyond baseline priors. When model probabilities are benchmarked against market odds to simulate historical value betting, the strategy yields positive returns on investment (ROI) and effectively mitigates maximum drawdown. The framework and data pipeline established in this project provide a robust foundation for multi-dimensional sports analytics and quantitative decision-making in sports betting markets.